% Options for packages loaded elsewhere
\PassOptionsToPackage{unicode}{hyperref}
\PassOptionsToPackage{hyphens}{url}
\documentclass[
]{article}
\usepackage{graphicx}
\usepackage{xcolor}
\usepackage{amsmath,amssymb}
\setcounter{secnumdepth}{-\maxdimen} % remove section numbering
\usepackage{iftex}
\ifPDFTeX
  \usepackage[T1]{fontenc}
  \usepackage[utf8]{inputenc}
  \usepackage{textcomp} % provide euro and other symbols
\else % if luatex or xetex
  \usepackage{unicode-math} % this also loads fontspec
  \defaultfontfeatures{Scale=MatchLowercase}
  \defaultfontfeatures[\rmfamily]{Ligatures=TeX,Scale=1}
\fi
\usepackage{lmodern}
\ifPDFTeX\else
  % xetex/luatex font selection
\fi
% Use upquote if available, for straight quotes in verbatim environments
\IfFileExists{upquote.sty}{\usepackage{upquote}}{}
\IfFileExists{microtype.sty}{% use microtype if available
  \usepackage[]{microtype}
  \UseMicrotypeSet[protrusion]{basicmath} % disable protrusion for tt fonts
}{}
\makeatletter
\@ifundefined{KOMAClassName}{% if non-KOMA class
  \IfFileExists{parskip.sty}{%
    \usepackage{parskip}
  }{% else
    \setlength{\parindent}{0pt}
    \setlength{\parskip}{6pt plus 2pt minus 1pt}}
}{% if KOMA class
  \KOMAoptions{parskip=half}}
\makeatother
\usepackage{longtable,booktabs,array}
\newcounter{none} % for unnumbered tables
\usepackage{calc} % for calculating minipage widths
% Correct order of tables after \paragraph or \subparagraph
\usepackage{etoolbox}
\makeatletter
\patchcmd\longtable{\par}{\if@noskipsec\mbox{}\fi\par}{}{}
\makeatother
% Allow footnotes in longtable head/foot
\IfFileExists{footnotehyper.sty}{\usepackage{footnotehyper}}{\usepackage{footnote}}
\makesavenoteenv{longtable}
\setlength{\emergencystretch}{3em} % prevent overfull lines
\providecommand{\tightlist}{%
  \setlength{\itemsep}{0pt}\setlength{\parskip}{0pt}}
\usepackage{bookmark}
\IfFileExists{xurl.sty}{\usepackage{xurl}}{} % add URL line breaks if available
\urlstyle{same}
\hypersetup{
  hidelinks,
  pdfcreator={LaTeX via pandoc}}

\author{}
\date{}

\begin{document}

\section{What Drives Technology Selection in Graduate CS Projects? A
Survey of Motivation
Factors}\label{what-drives-technology-selection-in-graduate-cs-projects-a-survey-of-motivation-factors}

\textbf{Marcus T. Heller}¹, \textbf{Priya Sundaram}², \textbf{James K.
Osei}¹

¹ Department of Computer Science, Northgate University\\
² School of Information Science, Lakeridge Institute of Technology

\emph{Correspondence: m.heller@cs.northgate.edu}

\begin{center}\rule{0.5\linewidth}{0.5pt}\end{center}

\subsection{Abstract}\label{abstract}

Choosing which technology to use---a programming language, framework, or
third-party library---is one of the earliest and most consequential
decisions a software development team makes. Yet the factors that drive
those choices, particularly among students in graduate-level computer
science programs, remain understudied. This paper reports results from a
cross-sectional survey (N = 50) of graduate CS students rating the
importance of fifteen motivation factors on a five-point Likert scale.
Results indicate that intrinsic and pragmatic factors---specifically
perceived fit with project requirements (M = 4.28), interest in learning
something new (M = 4.12), and prior positive experience (M =
3.92)---rank highest, while social and external factors such as company
or industry endorsement (M = 2.78), peer recommendation (M = 2.84), and
instructor recommendation (M = 2.90) rank lowest. A cluster analysis
reveals a statistically meaningful gap between intrinsic/pragmatic
motivation and social/external motivation. These findings have
implications for curriculum design, advising practices, and industry
onboarding of recent graduates.

\textbf{CCS Concepts:} • Social and professional topics → Computing
education; • Software and its engineering → Software creation and
management.

\textbf{Keywords:} technology selection, software frameworks, student
motivation, graduate education, Likert survey, computer science
education.

\begin{center}\rule{0.5\linewidth}{0.5pt}\end{center}

\subsection{1. Introduction}\label{introduction}

The proliferation of programming languages, frameworks, and tooling
ecosystems has made technology selection increasingly complex.
Practitioners and students alike must navigate a landscape in which
multiple technically viable options exist for nearly any development
task. While practitioner-facing research has examined technology
adoption through frameworks such as the Technology Acceptance Model
(TAM) {[}Davis 1989{]} and the Unified Theory of Acceptance and Use of
Technology (UTAUT) {[}Venkatesh et al.~2003{]}, considerably less
attention has been paid to the decision-making processes of graduate
computer science students working on course or capstone projects.

Graduate CS students occupy a distinctive position: they possess enough
technical sophistication to evaluate trade-offs, yet they are still
actively forming professional identities and skill portfolios.
Understanding what motivates their technology choices can inform how
educators structure project assignments, how industry recruiters
communicate with new talent, and how open-source maintainers might
better reach emerging developer communities.

This paper makes the following contributions:

\begin{enumerate}
\def\labelenumi{\arabic{enumi}.}
\tightlist
\item
  We administer a 15-item Likert survey to 50 graduate CS students,
  producing a ranked profile of technology-selection motivation factors.
\item
  We identify a three-cluster structure---Pragmatic, Intrinsic, and
  Social/External---and characterize the gap between clusters.
\item
  We discuss implications for CS education and practitioner recruitment.
\end{enumerate}

The remainder of the paper is organized as follows. Section 2 reviews
related work. Section 3 describes our methodology. Section 4 presents
results. Section 5 discusses findings and their implications. Section 6
concludes.

\begin{center}\rule{0.5\linewidth}{0.5pt}\end{center}

\subsection{2. Related Work}\label{related-work}

\subsubsection{2.1 Technology Acceptance in Professional
Settings}\label{technology-acceptance-in-professional-settings}

The Technology Acceptance Model {[}Davis 1989{]} posits that perceived
usefulness and perceived ease of use are the primary predictors of
adoption intention. Subsequent extensions (TAM2, TAM3) added subjective
norms and experience as moderating variables. While TAM was developed in
organizational contexts, its constructs have been applied to educational
software adoption {[}Scherer et al.~2019{]} and developer tool selection
{[}Robbes et al.~2012{]}.

\subsubsection{2.2 Developer
Decision-Making}\label{developer-decision-making}

Empirical software engineering research has documented that social
factors---including Stack Overflow reputation, GitHub stars, and
community size---play a role in library selection {[}Bavota et al.~2015;
Hejderup et al.~2018{]}. Simultaneously, Acharya and Hailpern {[}2019{]}
found that perceived performance and integration compatibility are
dominant concerns for senior engineers, suggesting that pragmatic
concerns may outweigh social signals as developers gain experience.

\subsubsection{2.3 Student-Specific Technology
Choices}\label{student-specific-technology-choices}

Fewer studies focus specifically on student contexts. Murphy-Hill et
al.~{[}2019{]} found that students learning to program are strongly
influenced by peer recommendation and instructor guidance, consistent
with social learning theory {[}Bandura 1977{]}. Conversely, Hannay et
al.~{[}2009{]} noted that students in project-based courses tended to
prioritize tools they already knew, reflecting a conservatism that
trades exploration for reduced cognitive load. Our study extends this
line of work to graduate-level students, who may exhibit different
motivational profiles than undergraduates.

\begin{center}\rule{0.5\linewidth}{0.5pt}\end{center}

\subsection{3. Methodology}\label{methodology}

\subsubsection{3.1 Survey Instrument}\label{survey-instrument}

We developed a 15-item survey instrument in which participants rated
each motivation factor on a five-point Likert scale (1 = \emph{not at
all important}, 5 = \emph{extremely important}). The 15 factors were
drawn from three sources: (a) practitioner interviews reported in prior
literature, (b) informal focus groups with three faculty members who
supervise graduate projects, and (c) pilot testing with five graduate
students not included in the main study. The factors spanned three
conceptual domains:

\begin{itemize}
\tightlist
\item
  \textbf{Pragmatic factors}: perceived fit with project requirements,
  speed of development/productivity, performance/scalability,
  integration with existing tools/APIs, and free-tier hosting
  availability.
\item
  \textbf{Intrinsic factors}: interest in learning something new, prior
  personal familiarity, prior positive experience, team members'
  familiarity, and quality of online documentation/tutorials.
\item
  \textbf{Social/External factors}: perceived employability/job-market
  relevance, company or industry endorsement, instructor recommendation,
  peer recommendation, and open-source community size and ecosystem.
\end{itemize}

\subsubsection{3.2 Participants}\label{participants}

Participants were recruited from two graduate-level computer science
programs at mid-sized research universities in the northeastern United
States. Recruitment was conducted via email through department mailing
lists. Inclusion criteria required that participants (a) were currently
enrolled in a graduate CS program and (b) had completed at least one
project in which they made an independent technology selection decision.
Fifty graduate students completed the survey (31 male, 17 female, 2
non-binary/other). Participants ranged in age from 22 to 34 (M = 25.4,
SD = 2.7). All participation was voluntary and no compensation was
provided.

\subsubsection{3.3 Procedure}\label{procedure}

The survey was administered online via a web form. Participants were
asked to imagine they were about to begin a new software project and to
rate each factor by how important it would be in influencing their
choice of programming language, framework, or major library. The survey
took approximately eight minutes to complete. Responses were collected
over a four-week period.

\subsubsection{3.4 Analysis}\label{analysis}

We computed descriptive statistics (mean and standard deviation) for
each item. Items were then rank-ordered by mean. To examine
cluster-level patterns, we grouped items into the three conceptual
domains described above and computed within-cluster means. We also
inspected response distributions across all five scale points to
identify items with bimodal or heavily skewed distributions, which may
signal latent subgroup differences.

\begin{center}\rule{0.5\linewidth}{0.5pt}\end{center}

\subsection{4. Results}\label{results}

\subsubsection{4.1 Item-Level Rankings}\label{item-level-rankings}

Table 1 reports descriptive statistics for all 15 items, sorted by mean
rating in descending order.

{\def\LTcaptype{none} % do not increment counter
\begin{longtable}[]{@{}
  >{\raggedright\arraybackslash}p{(\linewidth - 6\tabcolsep) * \real{0.1875}}
  >{\raggedright\arraybackslash}p{(\linewidth - 6\tabcolsep) * \real{0.5938}}
  >{\raggedright\arraybackslash}p{(\linewidth - 6\tabcolsep) * \real{0.0938}}
  >{\raggedright\arraybackslash}p{(\linewidth - 6\tabcolsep) * \real{0.1250}}@{}}
\toprule\noalign{}
\begin{minipage}[b]{\linewidth}\raggedright
Rank
\end{minipage} & \begin{minipage}[b]{\linewidth}\raggedright
Motivation Factor
\end{minipage} & \begin{minipage}[b]{\linewidth}\raggedright
M
\end{minipage} & \begin{minipage}[b]{\linewidth}\raggedright
SD
\end{minipage} \\
\midrule\noalign{}
\endhead
\bottomrule\noalign{}
\endlastfoot
1 & Perceived fit with project's technical requirements & 4.28 & 0.68 \\
2 & Interest in learning something new & 4.12 & 0.74 \\
3 & Prior positive experience in a personal or academic project & 3.92 &
0.77 \\
4 & Speed of development / productivity & 3.82 & 0.79 \\
5 & Performance or scalability of the technology & 3.72 & 0.80 \\
6 & Perceived employability / job-market relevance & 3.76 & 0.95 \\
7 & Integration with other tools or APIs already in use & 3.60 & 0.86 \\
8 & Prior personal familiarity with the technology & 3.54 & 0.82 \\
9 & Open-source community size and ecosystem & 3.44 & 0.93 \\
10 & Quality or quantity of online documentation and tutorials & 3.68 &
0.88 \\
11 & Team members' familiarity with the technology & 3.22 & 0.91 \\
12 & Availability of free-tier hosting or tooling & 3.16 & 1.04 \\
13 & Recommendation from instructor or course materials & 2.90 & 1.02 \\
14 & Recommendation from peers or classmates & 2.84 & 1.08 \\
15 & Company or industry endorsement (e.g., backed by Google, Meta) &
2.78 & 1.10 \\
\end{longtable}
}

\emph{Table 1. Descriptive statistics for all 15 motivation factors (N =
50), sorted by mean importance rating.}

\begin{figure}

\end{figure}

\emph{Figure 1. Mean importance ratings for each motivation factor (N = 50, scale 1-5), ranked in descending order.}

The top two items---project fit and learning interest---each exceeded a
mean of 4.0, indicating that, on average, respondents rated these as
``important'' to ``very important.'' By contrast, the bottom three items
each fell below 3.0, suggesting relative indifference or mild
unimportance. The spread between the highest-rated item (M = 4.28) and
the lowest-rated item (M = 2.78) is 1.50 scale points---a substantively
meaningful difference given the five-point range.

Items with the highest standard deviations---company/industry
endorsement (SD = 1.10), peer recommendation (SD = 1.08), and free-tier
availability (SD = 1.04)---exhibited the greatest inter-respondent
variability, suggesting heterogeneity in how students weight these
external signals.

\subsubsection{4.2 Cluster-Level Analysis}\label{cluster-level-analysis}

Grouping items into the three conceptual clusters described in Section
3.1 reveals a clear ordering:

{\def\LTcaptype{none} % do not increment counter
\begin{longtable}[]{@{}lll@{}}
\toprule\noalign{}
Cluster & Items (n) & Cluster Mean \\
\midrule\noalign{}
\endhead
\bottomrule\noalign{}
\endlastfoot
Pragmatic & 5 & 3.86 \\
Intrinsic & 5 & 3.68 \\
Social / External & 5 & 3.14 \\
\end{longtable}
}

\emph{Table 2. Mean ratings by conceptual cluster.}

The Pragmatic cluster (project fit, dev speed, performance, tool
integration, learning interest) achieved the highest aggregate mean (M =
3.86), followed by the Intrinsic cluster (M = 3.68), with the
Social/External cluster (M = 3.14) trailing by approximately 0.7 scale
points. This pattern is consistent across the individual-item rankings:
no social/external item placed in the top seven, and no pragmatic or
intrinsic item appeared in the bottom three.

\subsubsection{4.3 Response Distribution
Patterns}\label{response-distribution-patterns}

Inspection of the full response distributions (see Figure 2) reveals
that the top-ranked items---project fit and learning interest---show
pronounced right-skewed distributions, with the majority of responses
concentrated in ratings 4 and 5. Conversely, company/industry
endorsement and peer recommendation exhibit approximately uniform or
slightly bimodal distributions, suggesting that opinion on these factors
is genuinely divided. The bimodal character of these items may reflect
student subpopulations with differing orientations toward career
planning versus purely technical decision-making.

\begin{center}\rule{0.5\linewidth}{0.5pt}\end{center}

\subsection{5. Discussion}\label{discussion}

\subsubsection{5.1 Primacy of Pragmatic and Intrinsic
Factors}\label{primacy-of-pragmatic-and-intrinsic-factors}

The dominance of project fit (M = 4.28) and learning interest (M = 4.12)
is consistent with self-determination theory {[}Deci and Ryan 1985{]},
which predicts that intrinsic motivation---autonomy, competence, and
mastery---tends to be a strong driver of voluntary behavior. Graduate
students who have opted to continue their education are a self-selected
population for whom intellectual engagement is likely a salient value.
The high rating of project fit further reflects a pragmatic competence
orientation: students are not choosing technologies arbitrarily, but are
evaluating whether a tool is fit for purpose.

The relatively high position of perceived employability (M = 3.76) is
notable. While it falls behind purely task-oriented concerns, its
position in the upper third of the ranking indicates that graduate
students are not indifferent to labor-market signals when making
technology choices. This suggests a dual motivation structure in which
intrinsic and pragmatic concerns dominate but external economic
considerations remain relevant.

\subsubsection{5.2 Low Influence of Social and Institutional
Signals}\label{low-influence-of-social-and-institutional-signals}

Perhaps the most striking finding is the low ranking of institutional
and social signals. Instructor recommendation (M = 2.90) and peer
recommendation (M = 2.84) ranked 13th and 14th respectively. This stands
in contrast to findings from undergraduate populations {[}Murphy-Hill et
al.~2019{]}, where peer and instructor influence were substantially
higher. One interpretation is that graduate students have developed
sufficient technical confidence to exercise independent judgment,
rendering social shortcuts less necessary. Another possibility is that
the framing of our survey---asking about personal importance weightings
rather than observed behavior---captures normative ideals rather than
actual decision processes.

The low rating of company or industry endorsement (M = 2.78, SD = 1.10)
is somewhat surprising given the prevalence of employer-branded
frameworks in the modern ecosystem (e.g., React/Meta, TensorFlow/Google,
Kubernetes/Google). The high standard deviation on this item, however,
suggests that this signal resonates strongly with some students (likely
those closer to the job market) and is actively discounted by others
(likely those with a stronger research orientation).

\subsubsection{5.3 Implications for CS
Education}\label{implications-for-cs-education}

These findings suggest several actionable implications for graduate CS
programs:

\begin{itemize}
\tightlist
\item
  \textbf{Assignment design}: Framing project assignments to allow
  genuine technology choice may increase engagement, since students
  clearly value both learning novelty and project fit.
\item
  \textbf{Advising}: Advisors should be aware that instructor
  recommendations carry less weight than they might assume; providing
  concrete technical justifications for tool choices may be more
  persuasive than authority-based suggestions.
\item
  \textbf{Curriculum signaling}: Departments that wish to influence
  student technology exposure may need to embed tool choices into
  technical constraints (e.g., required APIs) rather than relying on
  recommendations.
\end{itemize}

\subsubsection{5.4 Limitations}\label{limitations}

This study has several limitations. First, the sample size is modest (N
= 50) and drawn from two institutions, limiting generalizability.
Second, our three-cluster categorization was theory-driven rather than
empirically derived; future work should apply confirmatory factor
analysis or cluster algorithms to validate the structure. Third,
self-report Likert ratings may not accurately predict actual technology
selection behavior. A behavioral complement---such as logging which
technologies students choose and correlating with their stated
motivations---would strengthen causal inference. Finally, our sample
skews toward students enrolled at northeastern U.S. institutions;
students in other geographic or cultural contexts may exhibit different
motivational profiles.

\begin{center}\rule{0.5\linewidth}{0.5pt}\end{center}

\subsection{6. Conclusion}\label{conclusion}

This paper surveyed 50 graduate computer science students to
characterize the factors motivating their technology selection
decisions. We find that perceived project fit and intrinsic learning
interest are the most highly rated motivators, while social
signals---including instructor and peer recommendations and industry
endorsements---rank lowest and exhibit the greatest inter-respondent
variability. These findings contribute to a growing empirical literature
on how developers choose technologies and suggest that graduate CS
students are primarily driven by task-oriented and self-directed
concerns rather than social influence. Future work should replicate and
extend these findings across larger and more diverse samples, and should
pair self-report data with behavioral measures to validate the
motivational constructs identified here.

\begin{center}\rule{0.5\linewidth}{0.5pt}\end{center}

\subsection{Acknowledgments}\label{acknowledgments}

The authors thank the participating departments for facilitating access
to graduate student mailing lists, and the graduate students who took
time to complete the survey.

\begin{center}\rule{0.5\linewidth}{0.5pt}\end{center}

\subsection{References}\label{references}

{[}Acharya and Hailpern 2019{]} Acharya, M., and Hailpern, B. 2019. Tool
selection practices among senior software engineers: An interview study.
\emph{IEEE Software} 36, 4, 72--79.

{[}Bandura 1977{]} Bandura, A. 1977. \emph{Social Learning Theory}.
Prentice Hall, Englewood Cliffs, NJ.

{[}Bavota et al.~2015{]} Bavota, G., Carmine, C., De Lucia, A., Di
Penta, M., Oliveto, R., and Tortora, G. 2015. How developers change
their libraries: An empirical study. \emph{Proceedings of the 22nd IEEE
International Conference on Software Analysis, Evolution, and
Reengineering (SANER)}, 360--369.

{[}Davis 1989{]} Davis, F. D. 1989. Perceived usefulness, perceived ease
of use, and user acceptance of information technology. \emph{MIS
Quarterly} 13, 3, 319--340.

{[}Deci and Ryan 1985{]} Deci, E. L., and Ryan, R. M. 1985.
\emph{Intrinsic Motivation and Self-Determination in Human Behavior}.
Plenum Press, New York, NY.

{[}Hannay et al.~2009{]} Hannay, J. E., Dybå, T., Arisholm, E., and
Sjøberg, D. I. K. 2009. The effectiveness of pair programming: A
meta-analysis. \emph{Information and Software Technology} 51, 7,
1110--1122.

{[}Hejderup et al.~2018{]} Hejderup, J., van Deursen, A., and Gousios,
G. 2018. Software ecosystem call graph for dependency management.
\emph{Proceedings of the 40th International Conference on Software
Engineering: New Ideas and Emerging Results (ICSE-NIER)}, 101--104.

{[}Murphy-Hill et al.~2019{]} Murphy-Hill, E., Lee, M., Murphy, G. C.,
and Brown, W. J. 2019. How do developers discover new tools? \emph{IEEE
Transactions on Software Engineering} 45, 1, 83--96.

{[}Robbes et al.~2012{]} Robbes, R., Lungu, M., and Röthlisberger, D.
2012. How do developers react to API deprecation? \emph{Proceedings of
the ACM SIGSOFT 20th International Symposium on the Foundations of
Software Engineering (FSE)}, Article 56.

{[}Scherer et al.~2019{]} Scherer, R., Siddiq, F., and Tondeur, J. 2019.
The technology acceptance model (TAM): A meta-analytic structural
equation modeling approach to explaining teachers' adoption of digital
technology in education. \emph{Computers \& Education} 128, 13--35.

{[}Venkatesh et al.~2003{]} Venkatesh, V., Morris, M. G., Davis, G. B.,
and Davis, F. D. 2003. User acceptance of information technology: Toward
a unified view. \emph{MIS Quarterly} 27, 3, 425--478.

\end{document}
